% Noether P09 Korean producer draft — UNCHECKED
% T04-U25; ED0002 lines6625--6648; source 1488 LF B / 6B4C76BEF6414656D9DA6E5F7AAED1014BB046D621D71B7387315DBD50A11131
% Pointer v009 / ED0002: B06BE3530D9CF2E82B56FDBA7FE41D5D044DF2425DFA2A059D4939EAA2F7A6C2 / C9A125167ACB33D914EE4374B65AE7CDF0052F568371B8B77B720EA178ABF0E3
% Translation only; proof fidelity remains checker-held.

이 모듈 영역이 성질 I--IV를 만족함은 쉽게 증명할 수 있다. 우선 꼴 (1)의 두 양을 더하거나 빼거나 곱해도 다시 꼴 (1)의 양이 되므로 I이 성립한다. 또한 $\mathfrak G$는 모든 기저의 수 $\eta$를 포함하고, $g(\eta)$가 $\eta$의 모든 정수적 대수함수를 두루 취할 때 모든 양 $\eta\cdot g(\eta)$도 포함한다. 따라서 $\eta$의 모든 \emph{정수적} 대수함수, 나아가 $\eta$의 모든 대수함수를 $\mathfrak G$의 두 양의 몫으로 나타낼 수 있으므로 II가 증명된다. 더 나아가 표현 (1)은 특히 $g_1=0\cdots g_t=0$으로 놓을 때 모든 대수적 정수를 포함한다. 그러나 $z$는 어떤 분수형 대수수와도 같아질 수 없다. 실제로 $z$가 대수적 수 $\beta$를 나타낸다면
\[
 \beta=\alpha+g_1(\eta)\eta_1+\cdots+g_t(\eta)\eta_t
\]
가 $\eta$에 관하여 항등적으로 성립하므로 반드시 $\beta=\alpha$이다. 이는 다음 특수화에서 알 수 있다.
\[
 \eta_1=0\cdots \eta_t=0.
\]
이때 승수 $g_1(\eta)\cdots g_t(\eta)$는 정수적 대수함수이므로 유한한 값을 유지한다. 곧 정수적 대수함수나 수로 바뀐다.\srcfn{*)}{IV에 대한 이 더 자세한 증명은 이 절 끝의 확장된 예 때문에 제시하였다. 여기서는 $z$가 구성상 정수적 대수함수라는 지적만으로도 충분하다.} 따라서 $\mathfrak G$는 조건 I--IV를 만족한다.
